%!TEX root = ED_Algoritmos.tex

\section{Introdução}
\begin{frame}<handout:0>[label=bcc]
	\frametitle{Plano Político-Pedagógico para o Bacharelado em Ciência da Computação}
	``A Formação Básica, como dito, tem por objetivo introduzir os conteúdos necessários ao desenvolvimento tecnológico e científico da Computação. A opção por disciplinas que compõem a Formação Básica este Plano leva em consideração as seguintes premissas:
	\setbeamercovered{transparent}
	\begin{enumerate}[<+->]
		\item O conceito de algoritmo é central para a Ciência da Computação;
		\item O conceito de algoritmo é intrinsecamente associado ao conceito de máquina;
		\item A habilidade de propor soluções algorítmicas corretas e eficientes é a essência
		da Ciência da Computação; e
		\item A habilidade de programar eficientemente é essencial ao desenvolvimento das
		atividades do cientista da Computação.''
	\end{enumerate}
\end{frame}

\begin{frame}<handout:0>[label=main]
	\framesubtitle{Como os algoritmos mudaram o mundo\footnote{\url{http://collegedegreesearch.net/algorithms/}}}

	Algoritmo: ``conjunto de instruções para resolver um problema ou realizar uma tarefa''.
	\\\vfill\pause
	\begin{columns}
		\column{.65\textwidth}
			\begin{itemize}[<+->]
			%Em matemática, o algoritmo de Euclides[a] é um método simples e eficiente de encontrar o máximo divisor comum entre dois números inteiros diferentes de zero. É um dos algoritmos mais antigos, conhecido desde que surgiu nos Livros VII e X da obra Elementos de Euclides[1] por volta de 300 a.C.. 
				\item Computadores
				\item Comunicações (TCP/IP, NTP, etc.)
				\item Busca e Localização (São Google!)
				\item Militares
				\item GPS, jogos digitais, etc.
			\end{itemize}
		\column{.3\textwidth}
			\includegraphics[width=\textwidth]{al-Khwarizmi}%
			\\\small\centering{al-Khwarizmi}
	\end{columns}
\end{frame}

\begin{frame}[label=main]
	%http://iconbug.com/detail/icon/2034/little-robot/
	\begin{center}
		
\includegraphics[height=10\baselineskip]{LittleRobotIcon}
	\end{center}
\end{frame}